\documentclass[11pt]{article}
\usepackage[T1]{fontenc}
\usepackage{lmodern,amsmath,amssymb,amsthm,booktabs,geometry,microtype}
\usepackage[colorlinks,allcolors=blue]{hyperref}
\geometry{left=0.94in,right=0.94in,top=0.85in,bottom=0.85in}
\newtheorem{theorem}{Theorem}
\newtheorem{lemma}[theorem]{Lemma}
\newtheorem{corollary}[theorem]{Corollary}
\newcommand{\Q}{\mathbb Q}
\newcommand{\F}{\mathbb F}
\newcommand{\Gal}{\operatorname{Gal}}
\newcommand{\GL}{\operatorname{GL}}
\newcommand{\PSL}{\operatorname{PSL}}
\newcommand{\ord}{\operatorname{ord}}
\newcommand{\A}{\mathcal A}
\setlength{\emergencystretch}{2em}
\raggedbottom
\title{Uniform Rasmussen--Tamagawa bounds over two families of sextic fields}
\author{}
\date{November 25, 2026}
\hypersetup{pdftitle={Uniform Rasmussen--Tamagawa bounds over two families of sextic fields}}
\begin{document}
\maketitle
\begin{abstract}
For every dimension $g$, there is one bound on the primes at which a
$g$-dimensional abelian variety over any $A_4$-sextic or $A_5$-sextic
field can be heavenly. The bound is independent of the discriminant, the
variety and its endomorphisms. For elliptic curves, $4{,}027$ suffices and only 29 primes pass the necessary tests.
The argument combines the odd residue degrees in two permutation
actions with the small-prime obstruction of Rasmussen and Tamagawa.
Dirichlet GRH supplies the entire analytic input, even when the
normal closure is nonsolvable.
\end{abstract}

\section{The uniformity question}

Call $A/K$ \emph{heavenly at $\ell$} if $K(A[\ell^\infty])$ lies in the
maximal pro-$\ell$ extension of $K(\mu_{\ell^\infty})$ unramified
outside $\ell$. Equivalently, $A$ has good reduction away from
$\ell$ and the mod-$\ell$ torsion extension over $K(\mu_\ell)$ has
$\ell$-power degree. Write $\A(K,g,\ell)$ for the set of
$K$-isomorphism classes of such varieties of dimension $g$.
Rasmussen and Tamagawa's uniform Conjecture~2 asks for emptiness for
$\ell>B(g,d)$ over all fields $K$ of degree $d$.
Their Theorem~5.1 assumes suitable Dedekind zeta hypotheses and treats
each fixed $K$. Their Theorem~5.2 proves uniformity in odd degree under
GRH. The comparison here uses the numbering of their author manuscript
\cite{rt}. Lombardo established uniformity in any degree for CM
abelian varieties \cite{lombardo}. McLeman and Rasmussen's later
quadratic work includes finiteness at fixed $\ell$ and results on
elliptic curves \cite{mcr}.

\begin{theorem}\label{main}
For every positive integer $g$ there is an effective $B(g)$ such that
$\A(K,g,\ell)=\varnothing$ whenever $[K:\Q]=6$, the normal closure
has group $A_4$ or $A_5$, and $\ell>B(g)$. One may take $B(1)=4{,}027$.
\end{theorem}
The two groups have unique transitive actions on six points up to
equivalence. For $A_4$ the stabilizer has order two and all such
subgroups are conjugate. In $A_5$ a subgroup of order ten has a unique
Sylow $5$-subgroup, hence is its normalizer. All these normalizers are
conjugate. Concrete fields occur at LMFDB labels
\href{https://www.lmfdb.org/NumberField/6.6.284124736.1}{6.6.284124736.1}
and \href{https://www.lmfdb.org/NumberField/6.2.27071209.2}{6.2.27071209.2}.

The earlier ProveIt report \cite{prior} already identifies these two
actions, proves the residue-degree argument including ramification, and
uses them to bound Momose type~2 isogeny primes. We apply that
arithmetic observation to a different obstruction, valid for every
dimension. The resulting theorem allows every endomorphism algebra and
includes primes dividing the discriminant of $K$. We found no earlier
statement of this uniform heavenly result for the two sextic families
in the sources reviewed. This is a paper proof using the upstream
Dirichlet nonvanishing theorem and published arithmetic inputs. The
downstream deductions have not been formalized in Lean.

\section{Two local inputs}

\begin{lemma}[Odd residue degree, following the earlier report]\label{cycles}
Suppose the normal closure $N/\Q$ of $K$ acts on the embeddings
$X=G/H$ of $K$, with $G=\Gal(N/\Q)$ and $H=\Gal(N/K)$.
If each element of $G$ has an odd cycle of length at most $r$ on $X$,
then each rational prime $p$ has a place $v$ of $K$ above it
of odd residue degree $f_v\le r$.
\end{lemma}
\begin{proof}
Choose a place of $N$ above $p$, its decomposition group $D$, its
inertia group $I$, and a lift $u\in D$ of residue-field Frobenius.
Places of $K$ above $p$ correspond to $D$-orbits on $X$.
In the orbit of $x$, with stabilizer $D_x$, the inertia orbits are
permuted transitively by $u$. Their number is $[D:I D_x]=f_v$.
Indeed the local degrees for $N_w/\Q_p$ and its subfield fixed by
$D_x$ give this index upon dividing the degree $[D:D_x]$ by its
ramification index $[I:I\cap D_x]$.
Thus the length of the $u$-cycle through $x$ is a multiple of $f_v$.
Choose $x$ on the asserted odd cycle. This argument also applies when
$I$ is nontrivial.
\end{proof}

\begin{lemma}[Rasmussen--Tamagawa's obstruction]\label{RT}
For fixed $g,d$ there exist an effective threshold $C_7(g,d)$
and a finite effective set of integers $\mathcal M(g,d)$ with the
following property. If $[K:\Q]=d$, $A$ is heavenly at
$\ell>C_7(g,d)$, then there is $m\in\mathcal M(g,d)$ dividing
$\ell-1$ such that for every $v\mid p\ne\ell$,
\[
  f_v\ \text{odd},\qquad p^{f_v}<\ell/(4g)
   \quad\Longrightarrow\quad
   p^{f_v}\bmod\ell\ \notin(\F_\ell^\times)^m .
\]
\end{lemma}
\begin{proof}
Use \cite[Cor.~3.2, (3.4), Lem.~3.5, \S4.1, Prop.~4.2]{rt}.
Their hypothesis (A1) is heavenliness. The displayed $C_7$ in
(3.4) suffices for (A2) uniformly in the degree. If $e$ is the
gcd of their local indices $e_\lambda$, then
$m=m_{\Q}\mid e/2$ and $e\mid M'(2g)d$. Thus the positive divisors
of $M'(2g)d$ give a finite set. Here $M'(a)$ is the integer in their
Lemma~3.1, the gcd of the orders of $\GL_a(\F_{\ell'})$ over odd
primes $\ell'\ne\ell$. Their valuation formula attains each minimum at infinitely many odd primes, so removing $\ell$ preserves the gcd. For $p\ne\ell$ the cyclotomic character
on arithmetic Frobenius at $v$ is $p^{f_v}\bmod\ell$.

To see the sign of the obstruction, triviality modulo $m$-th powers
would make all of the residual characters
$\chi^{i_j+i_k-1}$ trivial on this Frobenius, where
$\chi^{i_1},\dots,\chi^{i_{2g}}$ are the diagonal characters of
$A[\ell]$. If $Q=p^{f_v}$ and $a_j$ is the integral trace of the
$j$-th power of Frobenius, this gives $a_2\equiv 2gQ\pmod\ell$.
Weil's bound $|a_2|\le2gQ$ and $4gQ<\ell$ force equality.
Every Frobenius eigenvalue is therefore $+\sqrt Q$ or $-\sqrt Q$.
Since $Q$ is a nonsquare and $a_1$ an integer, $a_1=0$.
The same residual characters give $a_1^2\equiv4g^2Q\pmod\ell$,
a contradiction. No condition on ramification of $p$ or $\ell$ in
$K$ was used.
\end{proof}

\section{Uniform transfer and the two actions}

\begin{theorem}\label{general}
Fix $g,d,r\ge1$. Over all fields $K$ of degree $d$ satisfying the
cycle hypothesis of Lemma~\ref{cycles}, a single effective bound
$B(g,d,r)$ excludes heavenliness for larger $\ell$.
\end{theorem}
\begin{proof}
Take the maximum of $C_7(g,d)$ and the constants
$C_6(m,g,r,\Q)$ from \cite[Prop.~2.7]{rt} over
$m\in\mathcal M(g,d)$. For $m\mid\ell-1$, let $L_{m,\ell}$
be the unique degree-$m$ subfield of $\Q(\zeta_\ell)$.
The only zeta hypothesis of that proposition here is RH for
$\zeta_{L_{m,\ell}}$. Its factorization into primitive Dirichlet
$L$-functions supplies it from the five public Lean exports
\cite{lean}; the zeta factor is covered by the ordinary RH export.
Apply the proposition to the identity automorphism. It yields
a rational prime $p\ne\ell$ with
\[
               p<(\ell/(4g))^{1/r},\qquad
               p\bmod\ell\in(\F_\ell^\times)^m .
\]
The residue condition expresses complete splitting in $L_{m,\ell}$.
For $m=1$ one may simply take $p=2$ for large enough $\ell$.
Lemma~\ref{cycles} supplies $v$ with odd $f_v\le r$,
so $p^{f_v}<\ell/(4g)$ and $p^{f_v}$ is still an $m$-th power.
Lemma~\ref{RT} gives the contradiction. All fields used for the
analytic estimate are cyclotomic subfields, with bounded degree.
\end{proof}

For the sextic actions, the counts of elements of each cycle shape are
\[
\begin{array}{c|rrrr}
 &1^6&1^2 2^2&3^2&1\,5\\ \hline
 A_4&1&3&8&0\\
 A_5&1&15&20&24
\end{array}
\]
The three double transpositions in $A_4$ each fix two edges of a
tetrahedron. A $3$-cycle has two orbits of length three on its edges.
For $A_5$, each of the six point stabilizers has order ten and
contains five involutions. The fifteen involutions are conjugate, so
each fixes $6\cdot5/15=2$ points. An element of order five fixes its
unique Sylow subgroup, and one of order three fixes none.
Thus $r=3$ works for both families. Theorem~\ref{general} proves
the first assertion of Theorem~\ref{main}. The same argument on
inertia orbits explains why the field discriminant never enters
the answer. The companion script exhausts the sixteen transitive
sextic actions. These two are exactly the ones in which \emph{every}
element has an odd cycle. This is a classification of the sufficient
group condition. Other actions require an additional argument.

\section{An explicit bound in dimension one}

Let $E/K$ be a heavenly elliptic curve in either family. First take
$\ell>5$. Its mod-$\ell$ representation on $G_K$ has semisimple
characters $\chi^i,\chi^{1-i}$ for some integer $0\le i<\ell-1$
by \cite[Lem.~3.3]{rt} and the Weil pairing. These exponents describe
the characters on $G_K$ even if the image of $\chi|_{G_K}$ is proper.
Since $\ell\nmid|G|$, the inertia group above $\ell$ in $N$ is tame
and cyclic. The lengths of the cycles of an inertia generator on the
embeddings equal the ramification indices on their inertia orbits.
The cycle table therefore gives a place $\lambda\mid\ell$ with
$e(\lambda/\ell)=1$ or $3$. This use of inertia at the torsion prime
is separate from the Frobenius argument at auxiliary primes.

At any $\lambda\mid\ell$, an elliptic curve
acquires semistable reduction over an extension of ramification index
at most six. For completeness, take a short integral equation
$y^2=x^3+ax+b$ over the local field, with discriminant $\Delta$.
Write $v(0)=+\infty$. If $v(j)\ge0$, then $s=v(\Delta)\le3v(a)$.
From $\Delta=-16(4a^3+27b^2)$ it follows that
$s=3v(a)$ or $s=2v(b)$, since $\ell>3$.
Also $s\le2v(b)$. Scale the equation by a parameter of valuation
$s/12$. Such a parameter exists in a totally ramified extension of
degree $12/\gcd(s,12)\le6$. The new coefficients are integral and
the new discriminant a unit. If $v(j)<0$, then
$3v(a)=2v(b)<v(\Delta)$. Write $v(a)=2t$ and $v(b)=3t$.
Scaling by a parameter of valuation $t/2$ yields integral unit
coefficients, unit $c_4$ and discriminant of positive valuation.
The resulting reduction is multiplicative. The case $j=0$ belongs
to the first alternative.

In the good-reduction case the index just exhibited is a divisor of
$12$ other than $12$. In the multiplicative case it divides two.
The minimal semistability index divides any successful index by
\cite[\S3.1]{rt}. At our chosen place we consequently have
\[
 e=e_{E,K_\lambda}\,e(\lambda/\ell)
       \in\{1,2,3,4,6,9,12,18\}.
\]
The Tate--Oort relation in \cite[(3.1)]{rt}, valid before their
large-prime assumption (A2), supplies weights $0\le j,j'\le e$
with $ei\equiv j$ and $e(1-i)\equiv j'\pmod{\ell-1}$.
For $\ell>211$ their sum must equal $e$, since $2e<\ell-1$.
Interchanging the two characters permits $j\le e/2$.

\paragraph{A finite test that forces equal inertia weights.}
O'Hara uses powers of auxiliary Frobenius traces to test local
weights over cubic fields \cite[\S2.3]{ohara}. We use the same
trace identity, retaining both possible odd residue degrees and
combining the products below by a gcd.
Fix a prime $p\in\{2,3,5,7,11,13\}$. The residue field at some
$v\mid p$ has size $Q=p$ or $p^3$. At this place $E$ has good
reduction. Let $a$ be its Frobenius trace, with $a^2\le4Q$.
The traces of powers satisfy
\[
 S_0=2,\quad S_1=a,\quad S_k=aS_{k-1}-Q S_{k-2}.
\]
Reduction of the two eigenvalues mod $\ell$ and the two inertia
congruences give
\[
   S_e(a,Q)\equiv Q^j+Q^{e-j}\pmod\ell.
\]
For $j<e/2$ the positive right-hand integer is strictly larger
than $2Q^{e/2}$ by strict arithmetic--geometric mean, whereas
$|S_e|\le2Q^{e/2}$. Define
\[
 D_{e,j}=\gcd_{p\in\{2,3,5,7,11,13\}}
     \ \prod_{Q\in\{p,p^3\}}
     \prod_{a^2\le4Q}\bigl(Q^j+Q^{e-j}-S_e(a,Q)\bigr),
\]
where $a$ runs over rational integers. For unbalanced weights $\ell$
must divide this positive gcd. The largest prime divisor across
$0\le j<e/2$, for each possible $e$, is
\[
 \begin{array}{c|rrrrrrrr}
 e&1&2&3&4&6&9&12&18\\ \hline
 \text{largest prime}&13&13&109&37&109&211&109&211 .
 \end{array}
\]
The exact checker multiplies the positive terms and uses Euclid's
algorithm. It divides each gcd by primes through $211$ to exhaustion
and checks that the remainder equals one. This proves $j=e/2$ for
every candidate $\ell>211$.

\paragraph{What balance forces at a good prime.}
Now $(e/2)(2i-1)\equiv0\pmod{\ell-1}$. Hence $4\mid e$, as
$2i-1$ is odd, leaving $e=4$ or $12$. It also follows that
$\ell\equiv3\pmod4$ and
\[
                    6(2i-1)\equiv0\pmod{\ell-1}.
\]
If $p\bmod\ell$ is a nonzero square, so is $Q=p^f$. Therefore
$u=Q^{2i-1}$ has order dividing three: write $Q=b^2$ and use the
last congruence. The residual trace square is
\[
      a^2\equiv Q(2+u+u^{-1})\equiv Q\ \text{or }4Q\pmod\ell.
\]
If $4p^3<\ell$, then $0\le a^2\le4Q<\ell$, forcing equality
over the integers. Both $Q$ and $4Q$ are nonsquares since $f=1$ or
$3$. Consequently a candidate $\ell>211$ must have \emph{no} prime
$p$ with $4p^3<\ell$ that is a quadratic residue mod $\ell$.

\paragraph{Covering all large torsion primes.}
For $\ell\equiv3\pmod4$, quadratic reciprocity identifies such
residues with rational primes splitting in $\Q(\sqrt{-\ell})$,
including the supplementary law at two. Its absolute discriminant
is $\ell$. Bach and Sorenson \cite[Thm.~5.1]{bs} bound the least split
prime by $(4\log\ell+10)^2$ under the relevant quadratic zeta
hypothesis, available from Dirichlet GRH. Put $H=3\cdot10^{13}$.
For $y\ge\log H$, $4(4y+10)^6e^{-y}$ decreases because
$24/(4y+10)<1$. At $y=\log H<32$ its numerator is below
$4\cdot138^6<H$. Thus every $\ell>H$ has a forbidden prime.

From $253{,}507$ to $H$ we replayed the earlier report's
\texttt{cubic\_tail.cpp} extending only its upper endpoint.
It checks the stronger strict split-prime condition with two
additional nondivisibilities. Through $4\cdot127^3$ it tests
264,911 primes directly. Above this, the eligible primes
$2,3,5,7,11,13,17$ restrict the unsieved candidates to 1,440 classes
modulo $2{,}042{,}040$. The nonsquare and nonzero residue conditions
at primes $19$ to $127$ then give periodic bit masks on these
progressions. The code handles prime divisors of the extra
polynomials separately, so its reductions preserve every exception.
Of 782 integers surviving the masks, 246 are prime; each has a
witness. The eight higher exceptional primes were also checked.
The code and rerun transcript accompany this paper. Trial of
$4p^3<\ell$ directly through $253{,}507$ leaves 52 candidates
larger than 211, listed in \texttt{check\_elliptic.txt}.

\paragraph{Keeping the same diagonal character at every prime.}
For each remaining $\ell>211$ retain the exponents $0\le i<\ell-1$
with $6(2i-1)\equiv0\pmod{\ell-1}$. For $\ell\le211$ begin with
all exponents, imposing only for $\ell>5$ the necessary condition
that $ei\bmod(\ell-1)\le e$ for at least one admissible $e$.
For each prime $p\ne\ell$, discard $i$ if for \emph{both}
$Q=p,p^3$ the least absolute residue of
$Q^i+Q^{1-i}\pmod\ell$ exceeds $\lfloor2\sqrt{Q}\rfloor$.
For an actual curve its fixed exponent must survive every prime.
We may stop as soon as $\lfloor2\sqrt{p^3}\rfloor\ge\lfloor
\ell/2\rfloor$, since the $p^3$ option alone allows every residue
from then on. The Python checker yields the following necessary list:
\[
\begin{gathered}
 2,3,5,7,11,13,19,23,29,31,37,43,59,67,83,107,139,163,211,\\
 283,307,331,379,499,547,643,883,907,4027 .
\end{gathered}
\]
This completes the proof of the numerical part of
Theorem~\ref{main}. The tests permit the locally chosen places to
vary with $p$ and impose only necessary congruences. Existence of a
heavenly elliptic curve at a listed prime, especially at $4027$,
requires additional evidence.

\paragraph{Corollary for cubic fields.}
The same 29-prime necessary list holds for heavenly elliptic curves
over \emph{every} cubic number field. At any rational prime the
identity $\sum e_vf_v=3$ guarantees a place with odd $f_v=1$ or $3$.
If $\ell>5$, inertia in the normal closure is cyclic since its group
embeds in $S_3$, and its action on three points has an orbit of size
one or three. Thus the two local inputs to every numerical test
are identical. This supplies a numerical version of the uniform
cubic elliptic finiteness in \cite[Prop.~7.5]{rt}. For related
balance criteria and computational sieves in this very family,
see O'Hara \cite{ohara}.

\section{An infinite family of permutation actions}

Let $q\equiv1\pmod4$ be an odd prime power. In the natural action
of $\PSL_2(\F_q)$ on its $q+1$ projective lines, an element whose
eigenline belongs to $\mathbb P^1(\F_q)$ fixes a point.
In the other case, a determinant-one lift has eigenvalues
$\alpha,\alpha^q=\alpha^{-1}$ in $\F_{q^2}$. Therefore
$\alpha^{(q+1)/2}=\pm1$ and the corresponding projective element
has order dividing the odd integer $(q+1)/2$. Every one of its
cycles then has odd length at most $(q+1)/2$.
Theorem~\ref{general} proves the uniform heavenly bound in this
family for any fixed $g,q$, over every number field realizing this
action. The $q=5$ case recovers the $A_5$ sextic theorem. A further
extension follows from the tower identity $\sum_{w\mid v} e_wf_w=n$.
If a family already supplies an odd absolute residue degree at most
$r$ above every rational prime, any tower of fixed odd relative degree
$n$ supplies one at most $rn$. The proof of Theorem~\ref{general}
therefore applies uniformly to these towers as well.

\begin{thebibliography}{9}
\bibitem{rt} C. Rasmussen and A. Tamagawa,
\emph{Arithmetic of abelian varieties with constrained torsion},
Trans. Amer. Math. Soc. 369 (2017), 2395--2424.
\href{https://arxiv.org/abs/1302.1477}{Author manuscript, arXiv:1302.1477v1}.
\href{https://doi.org/10.1090/tran/6790}{Journal article}.
\bibitem{prior} ProveIt, \emph{Arithmetic transfers}, section on
uniform isogeny bounds for sextic fields, file \texttt{isogenies.tex}.
\href{https://github.com/VladimirReshetnikov/ProveIt/blob/b65aa6229415bdb66d9ce52decc4e997793cd07d/docs/incoming/grh-arithmetic-transfers.zip}{Pinned archive}.
\bibitem{lean} L. Sernau and collaborators, \emph{GRH for Dirichlet
$L$-functions, in Lean 4}, based on the work of L. Alp\"oge with Claude.
\href{https://github.com/lsernau/grh-lean/blob/72fead102f663d5c7a407964a13dc0cf9d977207/grh/src/Main.lean}{Pinned public theorems}.
\bibitem{bs} E. Bach and J. Sorenson, \emph{Explicit bounds for primes
in residue classes}, Math. Comp. 65 (1996), 1717--1735.
\href{https://doi.org/10.1090/S0025-5718-96-00763-6}{Journal article}.
\bibitem{lombardo} D. Lombardo, \emph{On the uniform
Rasmussen--Tamagawa conjecture in the CM case}, Math. Res. Lett. 25
(2018). \href{https://arxiv.org/abs/1511.09019}{Author manuscript}.
\bibitem{mcr} C. McLeman and C. Rasmussen, \emph{Heavenly elliptic
curves over quadratic fields}, Res. Number Theory 12 (2026), 79.
\href{https://doi.org/10.1007/s40993-026-00757-8}{Journal article}.
\bibitem{ohara} S. O'Hara, \emph{Heavenly elliptic curves over cubic
number fields} (2026). \href{https://arxiv.org/abs/2609.35642}{Author manuscript}.
\end{thebibliography}
\end{document}
